\documentclass[11pt,a4paper]{article}
\usepackage{gtu-el-yazisi}
\begin{document}
\notbilgi{Elektronik I}{Uygulama dersi: BJT ve MOSFET devre analizi}{Elektronik1-Uygulama-dersi.pdf}

\sayfa{1}
\section{Uygulama}

\begin{center}
\begin{circuitikz}[american]
  \draw (0,2) node[ocirc, label={[font=\small]above:$B$}] {} to[short] (1.5,2) to[R, l_=$r_\pi$] (1.5,0) to[short] (6.0,0);
  \node[font=\small] at (1.95,1.5) {$+$};
  \node[font=\small] at (1.95,0.5) {$-$};
  \node at (2.45,1.0) {$v_\pi$};
  \draw (3.0,0) to[short] (3.0,-0.5) node[ocirc, label={[font=\small]below:$E$}] {};
  \draw (4.0,2) to[cI, l=$g_m v_\pi$] (4.0,0);
  \draw (6.0,2) to[R, l=$r_o$] (6.0,0);
  \draw (4.0,2) to[short] (6.0,2) to[short] (7.2,2) node[ocirc, label={[font=\small]right:$C$}] {};
\end{circuitikz}
\end{center}

\[ r_o = \frac{V_A}{I_C} \]

\begin{center}
\begin{circuitikz}[american]
  \draw (0.5,3.2) to[short] (5.0,3.2) to[R, l_=$r_o$] (5.0,0) to[short] (0.5,0);
  \draw (0.5,3.2) to[short] (0.5,2.6) to[cI] (0.5,1.4) to[R, l_=$r_e$] (0.5,0);
  \draw (-1,1.4) node[ocirc, label={[font=\small]above:$B$}] {} to[short] (0.5,1.4);
  \node[font=\small, right] at (1.2,2.0) {$g_m v_\pi/\beta i_B/\alpha i_e$};
  \node[font=\small, right] at (1.2,0.7) {$\downarrow i_e$\quad $+\ v_\pi\ -$};
  \node[right] at (5.1,3.2) {$C$};
\end{circuitikz}
\end{center}

\elcizimi[0.4]{sekiller/Elektronik1-Uygulama-dersi-s1-1.png}[1]

\begin{ornek}
  Analyze the circuit at left. Assume $\beta = 100$, $V_A = \qty{100}{\volt}$. Calculate the voltage gain $\left(\dfrac{v_{out}}{v_{in}}\right)$, $R_{in}$, and $R_{out}$
  \begin{itemize}
    \item Determine the type of amplifier
      \begin{itemize}
        \item Common-Emitter (Use $\pi$ model)
      \end{itemize}
    \item DC analysis (to find small signal param\okunamadi{eters})
  \end{itemize}
\end{ornek}

$\Downarrow$ @DC

\elcizimi[0.4]{sekiller/Elektronik1-Uygulama-dersi-s1-2.png}[1]

\begin{align*}
  -\qty{3}{\volt} + \qty{100}{\kilo\ohm}\, I_B + 0{,}7 &= 0 \\
  \qty{100}{\kilo\ohm}\, I_B &= \qty{2,3}{\volt} \\
  I_B &= \qty{23}{\micro\ampere} \\
  I_E &= \qty{2,32}{\milli\ampere} \\
  I_C &= \qty{2,3}{\milli\ampere}
\end{align*}

\elcizimi[0.45]{sekiller/Elektronik1-Uygulama-dersi-s1-3.png}[1]

\[ V_{CB} \ge 0{,}4 \quad \text{(for linear)\footnotemark} \]
\footnotetext{Claude notu: muhtemelen $V_{CB} \ge -0{,}4\ \mathrm{V}$ olmalı.}
\begin{align*}
  r_o &= \frac{\qty{100}{\volt}}{\qty{2,3}{\milli\ampere}} \approx \qty{43,5}{\kilo\ohm} \\
  g_m &= \frac{I_C}{V_T} = \frac{\qty{2,3}{\milli\ampere}}{\qty{25}{\milli\volt}} = \qty{92}{\milli\siemens} \\
  r_e &= \frac{\okunamadi{26}\,\mathrm{mV}}{2{,}32} \approx \qty{11,2}{\ohm}\footnotemark \\
  r_\pi &= \qty{1,13}{\kilo\ohm} \qquad = (\beta+1)\,r_e
\end{align*}
\footnotetext{Claude notu: pay 26 mV ise $26/2{,}32 \approx 11{,}2\ \Omega$ olur; 25 mV ile $\approx 10{,}8\ \Omega$ çıkardı.}

\sayfa{2}
\begin{center}
\begin{circuitikz}[american, scale=0.8, transform shape]
  \draw (0,0) node[ground]{} to[V, l_=$v_{in}$] (0,2)
        to[R, l=\qty{100}{\kilo\ohm}] (2.6,2) node[circ]{} node[above=2pt]{$B$} to[short] (3.6,2)
        to[R, l_=\qty{1,13}{\kilo\ohm}] (3.6,0) to[short] (8.4,0);
  \node[font=\small] at (4.05,1.55) {$+$};
  \node[font=\small] at (4.05,0.45) {$-$};
  \node at (4.55,1.0) {$v_\pi$};
  \draw (5.4,0) to[short] (5.4,-0.4) node[ocirc, label={[font=\small]left:$E$}] {} to[short] (5.4,-0.8) node[ground]{};
  \draw (5.4,2) to[cI, l=\qty{92}{\milli\siemens}$\,v_\pi$] (5.4,0);
  \draw (8.4,2) to[R, l=\qty{43,5}{\kilo\ohm}] (8.4,0);
  \draw (5.4,2) to[short] (8.4,2) to[short] (10.8,2) node[ocirc, label={[font=\small]above:$C$}] {} to[short] (12.0,2)
        to[R, l=\qty{3}{\kilo\ohm}] (12.0,0) node[ground]{};
  \draw (12.0,2) to[short] (13.2,2) node[right]{$V_{out}$};
  \draw[-{Stealth[length=2mm]}] (-0.9,0.1) -- (-0.9,1.9) node[midway, left]{$R_{in}$};
\end{circuitikz}
\end{center}

\noindent
\begin{minipage}[t]{0.45\linewidth}
  \[ \begin{aligned} R_{in} &= 100 + 1{,}13 \\ &= \qty{101,13}{\kilo\ohm} \end{aligned} \]
\end{minipage}\hfill
\begin{minipage}[t]{0.45\linewidth}
  \[ \begin{aligned} R_{out} &= r_o \parallel R_C \\ R_{out} &\approx \qty{2,8}{\kilo\ohm} \end{aligned} \]
\end{minipage}

\begin{align*}
  V_{out} &= -\qty{2,8}{\kilo\ohm} \times \qty{92}{\milli\siemens}\; v_\pi \\
  v_\pi &= \frac{v_{in} \times 1{,}13\,\mathrm{k}}{\qty{101,3}{\kilo\ohm}}
\end{align*}
(ok: $v_\pi$ ifadesi $V_{out}$ denkleminin sonundaki $v_\pi$'ye yerleştiriliyor)
\[ V_{out} = \frac{-257{,}6 \cdot v_{in} \cdot 1{,}13\,\mathrm{k}}{101{,}3\,\mathrm{k}} = -2{,}68\ \mathrm{V/V}\footnotemark \]
\footnotetext{Claude notu: $257{,}6 \cdot 1{,}13/101{,}13 \approx 2{,}88$ olduğundan muhtemelen $-2{,}88\ \mathrm{V/V}$ olmalı.}

\elcizimi[0.5]{sekiller/Elektronik1-Uygulama-dersi-s2-1.png}[2]

\section{Soru 2}
\elcizimi[0.72]{sekiller/Elektronik1-Uygulama-dersi-s2-2.png}[2]

\noindent $\beta = 120$ \qquad $V_A = \infty$ \qquad C-E with $R_E$ (T model)

\sayfa{3}
\textbf{DC}
\[ -20 + 470\mathrm{k}\, I_B + 0{,}7 + 121 \cdot 0{,}56\mathrm{k}\, I_B = 0 \]
\[ \kutula{\begin{aligned} I_B &= \qty{0,035}{\milli\ampere} \\ I_C &= \qty{4,3}{\milli\ampere} \\ I_E &= \qty{4,34}{\milli\ampere} \end{aligned}} \]

\[ V_C = 20 - \qty{2,2}{\kilo\ohm} \times \qty{4,3}{\milli\ampere} = 10{,}54 \]

\begin{align*}
  r_e &= \frac{\qty{25}{\milli\volt}}{\qty{4,34}{\milli\ampere}} = \qty{5,76}{\ohm} \\
  g_m &= \frac{4{,}3}{25} = \qty{172}{\milli\siemens} \\
  r_o &= \infty
\end{align*}

\elcizimi[0.78]{sekiller/Elektronik1-Uygulama-dersi-s3-1.png}[3]

\[ R_{in} = \qty{470}{\kilo\ohm} \parallel \underbrace{(121)(\qty{565,76}{\ohm})}_{\qty{68,4}{\kilo\ohm}} \qquad\qquad R_{out} = \qty{2,2}{\kilo\ohm} \]
\[ R_{in} = \qty{59,8}{\kilo\ohm} \]

\begin{align*}
  v_{out} &= -\qty{172}{\milli\siemens} \cdot \qty{2,2}{\kilo\ohm}\, v_\pi \\
  v_\pi &= \frac{v_{in}}{565} \times 5{,}76 \\
  v_{out} &= -378{,}4\, v_{in} \times 0{,}01 \\
  \frac{v_{out}}{v_{in}} &\approx -3{,}784\ [\mathrm{V/V}]
\end{align*}

\sayfa{4}
\section{Soru}

$R_{in},\ R_{out},\ A_v\ \left(\dfrac{v_{out}}{v_{sig}}\right)$ \quad \textcolor{kirmizikalem}{CB Amplifier}

\elcizimi[0.7]{sekiller/Elektronik1-Uygulama-dersi-s4-1.png}[4]

\[ \beta = 100 \qquad V_A = \infty \]

\begin{align*}
  I_C &= \qty{1}{\milli\ampere} \\
  V_T &= \qty{25}{\milli\volt} \\
  r_e &\simeq \qty{25}{\ohm} \\
  r_\pi &\simeq \qty{2,5}{\kilo\ohm} \\
  g_m &= \frac{\qty{1}{\milli\ampere}}{\qty{25}{\milli\volt}} = \qty{40}{\milli\siemens}
\end{align*}

\elcizimi[0.55]{sekiller/Elektronik1-Uygulama-dersi-s4-2.png}[4]

\begin{align*}
  R_{in} &= r_e = \qty{25}{\ohm} \\
  R_{out} &= \qty{8}{\kilo\ohm}
\end{align*}

\textcolor{kirmizikalem}{Common emitter de sadece \okunamadi{faz} farkı var!}

\begin{align*}
  v_{out} &= -\qty{3,07}{\kilo\ohm}\,\alpha \cdot i_e \\
  i_e &= \frac{-v_{sig}}{\qty{5,025}{\kilo\ohm}} \\
  \frac{v_{out}}{v_{sig}} &= \frac{3{,}07 \cdot 0{,}99}{5{,}025} \approx 0{,}6\ [\mathrm{V/V}]
\end{align*}

\sayfa{5}
\section{Soru}

\elcizimi[0.9]{sekiller/Elektronik1-Uygulama-dersi-s5-1.png}[5]

\[ \beta = 100 \qquad V_A = \qty{100}{\volt} \]

\begin{align*}
  r_e &= \frac{\qty{25}{\milli\volt}}{\qty{5}{\milli\ampere}} = \qty{5}{\ohm} \\
  r_\pi &= \qty{0,5}{\kilo\ohm} \\
  g_m &= \frac{\qty{5}{\milli\ampere}}{\qty{25}{\milli\volt}} = \qty{200}{\milli\siemens} \\
  r_o &= \frac{\qty{100}{\volt}}{\qty{5}{\milli\ampere}} = \qty{20}{\kilo\ohm}
\end{align*}

\elcizimi[0.8]{sekiller/Elektronik1-Uygulama-dersi-s5-2.png}[5]

\begin{align*}
  R_{in} &= \qty{40}{\kilo\ohm} \parallel (\qty{0,5}{\kilo\ohm} + \qty{101}{\kilo\ohm}) \\
  R_{in} &\approx \qty{28,5}{\kilo\ohm} \\
  R_{out} &= \frac{(\qty{10}{\kilo\ohm} \parallel \qty{40}{\kilo\ohm}) + \qty{0,5}{\kilo\ohm}}{\beta+1} \\
  R_{out} &\approx \qty{80}{\ohm}
\end{align*}

\begin{align*}
  V_{out} &= \left( \frac{v_\pi}{0{,}5\mathrm{k}} + \qty{200}{\milli\siemens}\, v_\pi \right) 1\mathrm{k} \\
  &= 2\, v_\pi + 200\, v_\pi \\
  &= 202\, v_\pi
\end{align*}

\sayfa{6}
\noindent $R_{sig} = 0$ ise \quad $\Rightarrow\ v_{sig} = v_\pi + v_{out}$
\[ \frac{v_{out}}{v_{sig}} = \frac{202\, v_\pi}{203\, v_\pi} \approx 0{,}99\ [\mathrm{V/V}] \]

\noindent $R_{sig} = \qty{10}{\kilo\ohm}$

\begin{center}
\begin{circuitikz}[american]
  \draw (0,0) node[ground]{} to[sV, l_=$v_{sig}$] (0,2) to[R, l=\qty{10}{\kilo\ohm}] (3,2) node[circ]{} node[above]{$X$} to[R, l=\qty{28,5}{\kilo\ohm}] (3,0) node[ground]{};
\end{circuitikz}
\end{center}

\[ V_X = \frac{28{,}5}{38{,}5}\, v_{sig} \approx 0{,}74\, v_{sig} \]

\begin{align*}
  v_\pi &= V_X - v_{out} \\
  v_\pi &= 0{,}74\, v_{sig} - 202\, v_\pi \\
  203\, v_\pi &= 0{,}74\, v_{sig} \\
  v_{sig} &= 273{,}4\, v_\pi
\end{align*}
\[ \Rightarrow\quad \frac{v_{out}}{v_{sig}} = \frac{202\, v_\pi}{274{,}3\, v_\pi} \approx 0{,}734\ [\mathrm{V/V}] \]

\begin{ornek}
  Design the circuit given below to obtain a current of $\qty{80}{\micro\ampere}$. Find the value of $R_D$ and $V_D$. Let the NMOS have $V_T = \qty{0,6}{\volt}$, $\mu_n C_{ox} = \qty{200}{\micro\ampere\per\volt\squared}$, $L = \qty{0,8}{\micro\meter}$, $W = \qty{4}{\micro\meter}$. Neglect channel length modulation.
\end{ornek}

\elcizimi[0.4]{sekiller/Elektronik1-Uygulama-dersi-s6-1.png}[6]

\begin{align*}
  I_D &= \tfrac{1}{2}\,\mu_n C_{ox}\,\frac{W}{L}\,(V_{GS}-V_T)^2 \\
  \qty{80}{\micro\ampere} &= \tfrac{1}{2}\cdot\qty{200}{\micro\ampere\per\volt\squared}\cdot\frac{4}{0{,}8}\,(V_{ov})^2 \\
  \frac{160}{1000}\;(V_{ov})^2 &= 0{,}16
\end{align*}
\[ \swarrow\ 0{,}4 \qquad\qquad \searrow\ \textcolor{kirmizikalem}{-0{,}4}\ \text{(kırmızıyla çizilmiş)} \]

\sayfa{7}
\begin{align*}
  V_{GS} - V_T &= \qty{0,4}{\volt} \\
  V_{GS} &= \qty{0,4}{\volt} + \qty{0,6}{\volt} = \qty{1}{\volt} \\
  V_G &= \qty{1}{\volt} + V_S = \qty{1}{\volt} \\
  V_G &= V_D = \qty{1}{\volt} \\
  R_D &= \frac{3-1}{80\cdot 10^{-6}} \\
      &= \frac{2}{80\cdot 10^{-6}} = \frac{2\cdot 10^{6}}{80} = \qty{25}{\kilo\ohm}
\end{align*}

\[ \begin{aligned} &\text{Sat} \\ V_{GS} &\ge V_T \\ V_{DS} &\ge V_{GS} - V_T \end{aligned} \]

\elcizimi[0.7]{sekiller/Elektronik1-Uygulama-dersi-s7-1.png}[7]

\begin{align*}
  I_{D1} &= \qty{80}{\micro\ampere} \\
  V_{D1} &= ? \\
  I_{D2} &= ?
\end{align*}

\begin{align*}
  I_D &= \tfrac{1}{2}\,\mu_n C_{ox}\,\frac{W}{L}\,(V_{GS1}-V_T)^2 \\
  I_{D2} &= \tfrac{1}{2}\,\mu_n C_{ox}\,\frac{W}{L}\,(V_{GS2}-V_T)^2 \\
  I_{D1} &= I_{D2} = \qty{80}{\micro\ampere}
\end{align*}
\textit{(İlk iki satırda $\tfrac{1}{2}\mu_n C_{ox}\tfrac{W}{L}$ çarpanı çizilmiş; $V_{GS1}$ ve $V_{GS2}$ daire içine alınmış.)}

\[ R_{min} \le R \le R_{max} \]

\[ V_T = \qty{0,6}{\volt} \]
\[
\text{Sat} \left\{ \begin{aligned}
  & V_G \ge V_T \quad \text{Sat} & \frac{3-0{,}6}{\qty{80}{\micro\ampere}} &= \qty{30}{\kilo\ohm} \\
  & V_{DS} \ge V_{GS} - V_T & \frac{3-2}{\qty{80}{\micro\ampere}} &= \qty{12,5}{\kilo\ohm} \\
  & V_D = V_G - V_T & & \\
  & V_G = V_D + V_T = \qty{2}{\volt} & &
\end{aligned} \right.
\]

\[
\left.\begin{aligned} V_{DS} &\le V_{GS}-V_T \\ V_{GS} &\ge V_T \end{aligned}\right\}\ \text{triode}
\qquad
\left.\begin{aligned} V_{GS} &\ge V_T \\ V_{DS} &\ge \underbrace{V_{GS}-V_T}_{V_{ov}} \end{aligned}\right\}\ \text{SA}
\]
\[ V_{GS} < V_T \ \}\ \text{cutoff} \]

\sayfa{8}
\begin{ornek}
  Design the circuit to obtain a drain voltage of $\qty{0,1}{\volt}$. What is the effective resistance between drain and source at this point $V_T = \qty{1}{\volt}$ and $\tfrac{1}{2}\mu_n C_{ox}\tfrac{W}{L} = \qty{1}{\milli\ampere\per\volt\squared}$.
\end{ornek}

\elcizimi[0.38]{sekiller/Elektronik1-Uygulama-dersi-s8-1.png}[8]

\[ V_{DS} < V_{GS} - V_T \qquad\qquad 0{,}1 < \qty{5}{\volt} \ \Big\}\ \text{triode} \]
\textit{(İlk ifadenin altında dalgalı çizgi ve ``X'' işareti var.)}
\[ I_D = \mu_n C_{ox} \frac{W}{L} \left( V_{GS} - V_T - \frac{V_{DS}}{2} \right) V_{DS} \]
\[ \frac{V_{DS}}{I_D} = r_{DS} = r_x = \]
\[ \Rightarrow\ \frac{1000}{7{,}9} \approx \qty{125}{\ohm} \]

\noindent $R_D$'yi bul $= ?$

\section{Soru}
\elcizimi[0.8]{sekiller/Elektronik1-Uygulama-dersi-s8-2.png}[8]

\sayfa{9}
\begin{ornek}
  $|V_T| = \qty{2}{\volt}$, $\mu_n C_{ox} = \qty{125}{\micro\ampere\per\volt\squared}$, $W = \qty{4}{\micro\meter}$, $L = \qty{1}{\micro\meter}$
  \begin{itemize}
    \item Perform DC analysis, calculate small signal params
    \item Estimate the voltage gain $\left[\dfrac{v_{out}}{v_{sig}}\right]$, $R_{in}$, $R_{out}$
  \end{itemize}
  $V_A = \qty{50}{\volt}$
\end{ornek}

\begin{align*}
  \textcolor{kirmizikalem}{\text{DC:}}\quad I_D &= \tfrac{1}{2}\mu_n C_{ox}\tfrac{W}{L}\,(V_{GS}-V_T)^2 \\
  \textcolor{kirmizikalem}{\text{Small signal:}}\quad r_o &= V_A / I_D \\
  g_m &= 2 I_D / V_{ov}
\end{align*}

\elcizimi[0.5]{sekiller/Elektronik1-Uygulama-dersi-s9-1.png}[9]

\[ V_{TH} = \qty{5}{\volt} \qquad R_{TH} = \qty{50}{\kilo\ohm}\,\footnotemark \]
\footnotetext{Claude notu: $\qty{100}{\mega\ohm}\parallel\qty{100}{\mega\ohm}=\qty{50}{\mega\ohm}$ olduğundan muhtemelen $\qty{50}{\mega\ohm}$ olmalı.}

\elcizimi[0.55]{sekiller/Elektronik1-Uygulama-dersi-s9-2.png}[9]

\[ I_D = \tfrac{1}{2}\mu_n C_{ox}\tfrac{W}{L}\,(V_{GS}-V_T)^2 \]
\[
\left.\begin{aligned} V_G &= \qty{5}{\volt} \\ V_S &= I_D \end{aligned}\right\}
\qquad
\begin{aligned}
  I_D &= \tfrac{1}{2}\cdot 125\ \mu\mathrm{A/V^2}\cdot 4\,(5 - I_D - 2)^2 \\
  I_D &= 250\ \mu\mathrm{A/V^2}\,(5 - I_D - 2)^2
\end{aligned}
\]
\[ I_D = \qty{1}{\milli\ampere} \]

\begin{align*}
  r_o &= \frac{50}{1\,\mathrm{mA}} = \qty{50}{\kilo\ohm} \\
  g_m &= \frac{2 \times 1\,\mathrm{mA}}{2} \qquad\qquad g_m = \qty{1}{\milli\siemens}
\end{align*}

\sayfa{10}
\elcizimi[0.92]{sekiller/Elektronik1-Uygulama-dersi-s10-1.png}[10]

\[ R_{out} = 5 \parallel \qty{50}{\kilo\ohm} = \qty{4,54}{\kilo\ohm} \]

\begin{align*}
  v_{gs} &= \frac{v_{sig}\cdot \qty{50}{\mega\ohm}}{50{,}1} \approx v_{sig} \\
  v_{out} &= -\underbrace{(\qty{50}{\kilo\ohm} \parallel \qty{5}{\kilo\ohm} \parallel \qty{10}{\kilo\ohm})}_{-3{,}12\,\mathrm{k\Omega}} \cdot \qty{1}{\milli\siemens} \cdot \overset{v_{sig}}{v_{gs}} \\
  \frac{v_{out}}{v_{sig}} &= -3{,}12\ \mathrm{V}
\end{align*}

\section{Soru}
\elcizimi[0.65]{sekiller/Elektronik1-Uygulama-dersi-s10-2.png}[10]

\begin{ornek}
  $V_T = \qty{1,5}{\volt}$ \qquad $\mu_n C_{ox}\dfrac{W}{L} = \qty{1}{\milli\ampere\per\volt\squared}$ \qquad $|V_A| = \infty$
  \begin{enumerate}[label=\alph*)]
    \item DC analysis
    \item Draw small signal equivalent circuit
    \item Calculate gain, $R_{in}$, $R_{out}$
  \end{enumerate}
\end{ornek}

\noindent $I_D = \qty{0,5}{\milli\ampere}$ \qquad\qquad Common-Gate (T model)

\begin{align*}
  V_D &= \qty{2,5}{\volt} \\
  V_G &= \qty{0}{\volt} \\
  \qty{0,5}{\milli\ampere} &= \tfrac{1}{2}\,\qty{1}{\milli\ampere\per\volt\squared}\,(V_{ov})^2 \\
  V_{ov} &= \qty{1}{\volt} \\
  V_{GS} - V_T &= \qty{1}{\volt} \\
  0 - V_S - 1{,}5 &= \qty{1}{\volt} \qquad\qquad V_S = -\qty{2,5}{\volt}
\end{align*}

\sayfa{11}
\[ \begin{aligned} V_{DS} &\ge V_{ov} \\ \qty{5}{\volt} &\ge 1\ \mathrm{V} \quad \checkmark\ \text{Sat} \end{aligned} \]

\begin{align*}
  r_o &= \infty \\
  g_m &= \frac{2 I_D}{V_{ov}} = \frac{2 \cdot 0{,}5\ \mathrm{mA}}{1\ \mathrm{V}} \\
  g_m &= \qty{1}{\milli\siemens}
\end{align*}

\elcizimi[0.9]{sekiller/Elektronik1-Uygulama-dersi-s11-1.png}[11]

\begin{align*}
  v_{out} &= -\qty{7,5}{\kilo\ohm} \cdot \qty{1}{\milli\siemens}\, v_\pi \\
          &= -7{,}5\, v_\pi \\
  v_{sig} &\approx v_\pi \\
  \frac{v_{out}}{v_{sig}} &= -7{,}5\ [\mathrm{V/V}]
\end{align*}

\end{document}
